% Źródła: Eves s. 285--286, 289, 304; Wikipedia: Johann Heinrich Lambert, Lambert quadrilateral.
% https://en.wikipedia.org/wiki/Lambert_quadrilateral

\section{Johann Heinrich Lambert} % lang-pl
\section{Johann Heinrich Lambert} % lang-it
\label{sekcja_lambert}

Trzydzieści trzy lata po Saccherim niemiecki uczony pójdzie tą samą drogą, ale zajdzie dalej. % lang-pl
W tej sekcji zobaczymy, kim on będzie, jaką figurę wybierze za podstawową i jak nazwie to, czego nie potrafi dowieść. % lang-pl
Trentatré anni dopo Saccheri uno studioso tedesco percorrerà la stessa strada, ma andrà più lontano. % lang-it
In questa sezione vedremo chi sarà, quale figura sceglierà come fondamentale e come chiamerà ciò che non saprà dimostrare. % lang-it

Johann Heinrich Lambert urodzi się 26 lub 28 sierpnia 1728 roku w Miluzie, a w 1763 roku, dzięki poparciu Fryderyka II, wejdzie do Pruskiej Akademii Nauk w Berlinie, gdzie będzie pracował obok Eulera do śmierci 25 września 1777 roku. % lang-pl
W 1765 roku zacznie korespondować z Immanuelem Kantem, który zechce mu zadedykować \emph{Krytykę czystego rozumu} -- książka ukaże się jednak dopiero po śmierci Lamberta. % lang-pl
Johann Heinrich Lambert nascerà il 26 o 28 agosto 1728 a Mulhouse e nel 1763, grazie all'appoggio di Federico II, entrerà nell'Accademia delle Scienze di Prussia a Berlino, dove lavorerà accanto a Eulero fino alla morte, il 25 settembre 1777. % lang-it
Nel 1765 comincerà a scambiare lettere con Immanuel Kant, che vorrà dedicargli la \emph{Critica della ragion pura} -- il libro uscirà tuttavia solo dopo la morte di Lambert. % lang-it
\index[persons]{Lambert, Johann Heinrich}%
\index[persons]{Kant, Immanuel}%
\index[persons]{Euler, Leonhard}%

W 1766 roku Lambert napisze \emph{Die Theorie der Parallellinien}, która zostanie opublikowana dopiero po jego śmierci. % lang-pl
Za figurę podstawową weźmie czworokąt o trzech kątach prostych, czyli połowę czworokąta Saccheriego, i rozważy trzy hipotezy zależnie od tego, czy czwarty kąt jest ostry, prosty czy rozwarty \cite[s. 285--286]{eves1_1972}. % lang-pl
Nel 1766 Lambert scriverà \emph{Die Theorie der Parallellinien}, che sarà pubblicata solo dopo la sua morte. % lang-it
Come figura fondamentale prenderà un quadrilatero con tre angoli retti, cioè metà di un quadrilatero di Saccheri, e considererà tre ipotesi a seconda che il quarto angolo sia acuto, retto o ottuso \cite[p. 285--286]{eves1_1972}. % lang-it

\begin{definition}
[czworokąt Lamberta] % lang-pl
[quadrilatero di Lambert] % lang-it
	Czworokąt o trzech kątach prostych. % lang-pl
	Un quadrilatero con tre angoli retti. % lang-it
\index{czworokąt!Lamberta} % lang-pl
\index{quadrilatero!di Lambert} % lang-it
\end{definition}

Nazywa się go też czworokątem Ibn al-Hajsama--Lamberta, bo Ibn al-Hajsam rozważy go już w XI wieku. % lang-pl
Połączenie środków podstawy i górnej podstawy czworokąta Saccheriego dzieli go na dwa czworokąty Lamberta (Eves \cite[s. 289, zad. 9]{eves1_1972}). % lang-pl
Lo si chiama anche quadrilatero di Ibn al-Haytham--Lambert, perché Ibn al-Haytham lo considererà già nell'XI secolo. % lang-it
Il segmento che congiunge i punti medi della base e della base superiore di un quadrilatero di Saccheri lo divide in due quadrilateri di Lambert (Eves \cite[p. 289, es. 9]{eves1_1972}). % lang-it

Lambert pójdzie dalej niż Saccheri: pokaże, że w trzech hipotezach suma kątów trójkąta jest odpowiednio mniejsza niż, równa lub większa od dwóch kątów prostych, a ponadto że niedomiar (w hipotezie kąta ostrego) albo nadmiar (w hipotezie kąta rozwartego) jest proporcjonalny do pola trójkąta. % lang-pl
Zauważy, że hipoteza kąta rozwartego przypomina geometrię sferyczną, gdzie pole trójkąta jest proporcjonalne do nadmiaru kątów, i wypowie domysł, że hipotezę kąta ostrego dałoby się zrealizować na sferze o urojonym promieniu \cite[s. 286]{eves1_1972}. % lang-pl
Lambert andrà oltre Saccheri: mostrerà che nelle tre ipotesi la somma degli angoli di un triangolo è rispettivamente minore, uguale o maggiore di due angoli retti, e che inoltre il difetto (nell'ipotesi dell'angolo acuto) o l'eccesso (in quella dell'angolo ottuso) è proporzionale all'area del triangolo. % lang-it
Noterà che l'ipotesi dell'angolo ottuso ricorda la geometria sferica, dove l'area di un triangolo è proporzionale all'eccesso, e congetturerà che l'ipotesi dell'angolo acuto possa realizzarsi su una sfera di raggio immaginario \cite[p. 286]{eves1_1972}. % lang-it

Ten domysł podejmie Franz Taurinus (zobacz sekcję \ref{sekcja_schweikart_taurinus}), a Lambert, jako pierwszy, wprowadzi do matematyki funkcje hiperboliczne i obliczy pole trójkąta hiperbolicznego. % lang-pl
Hipotezę kąta rozwartego wyeliminuje tym samym cichym założeniem co Saccheri, a wnioski dotyczące kąta ostrego pozostaną niejednoznaczne i niezadowalające -- dlatego pracy nie opublikuje za życia \cite[s. 286]{eves1_1972}. % lang-pl
Questa congettura sarà ripresa da Franz Taurinus (si veda la sezione \ref{sekcja_schweikart_taurinus}), mentre Lambert introdurrà per primo in matematica le funzioni iperboliche e calcolerà l'area di un triangolo iperbolico. % lang-it
L'ipotesi dell'angolo ottuso sarà eliminata con la stessa tacita assunzione di Saccheri, mentre le conclusioni sull'angolo acuto resteranno vaghe e insoddisfacenti -- per questo il lavoro non sarà pubblicato in vita \cite[p. 286]{eves1_1972}. % lang-it
